% !TeX program = xelatex
% !TeX encoding = UTF-8
% Presentazione revisionata: Borealis: serie temporali ambientali
% Basata sul materiale originale fornito dall'autore e sul template M01_Pensare_con_i_dati.tex
% Compilazione: xelatex -interaction=nonstopmode Borealis_M09_serie_temporali.tex (due passaggi)
% Diagrammi TikZ e figure statistiche vettoriali; sorgenti R e dati inclusi.
% Revisione completa delle 106 slide; fonti consultate il 19 novembre 2026.
\documentclass[12pt,aspectratio=1610]{beamer}
\geometry{paperwidth=20cm,paperheight=12.5cm}
\usepackage[italian,provide=*]{babel}
\usepackage{fontspec}
\IfFontExistsTF{NimbusSans-Regular.otf}{
\setsansfont{NimbusSans-Regular.otf}[BoldFont=NimbusSans-Bold.otf,ItalicFont=NimbusSans-Italic.otf,BoldItalicFont=NimbusSans-BoldItalic.otf]
\setmainfont{NimbusSans-Regular.otf}[BoldFont=NimbusSans-Bold.otf,ItalicFont=NimbusSans-Italic.otf,BoldItalicFont=NimbusSans-BoldItalic.otf]
}{
\setsansfont{TeX Gyre Heros}
\setmainfont{TeX Gyre Heros}
}
\IfFontExistsTF{DejaVu Sans Mono}{
\setmonofont{DejaVu Sans Mono}[Scale=.88]
}{\setmonofont{Latin Modern Mono}[Scale=.88]}
\usefonttheme{professionalfonts}
\usepackage{amsmath,amssymb,booktabs,array,graphicx,tikz,bookmark}
\usetikzlibrary{calc,positioning,arrows.meta}
\graphicspath{{Risorse/Figure/}}
\definecolor{Paper}{HTML}{F8F6F1}
\definecolor{Ink}{HTML}{253344}
\definecolor{Navy}{HTML}{15283E}
\definecolor{Quantum}{HTML}{5146A5}
\definecolor{Classical}{HTML}{087C74}
\definecolor{Muted}{HTML}{657180}
\definecolor{Rule}{HTML}{D6DAD9}
\definecolor{Night}{HTML}{112236}
\definecolor{Lilac}{HTML}{ABB5FF}
\definecolor{Mint}{HTML}{72D8C8}
\setbeamercolor{background canvas}{bg=Paper}
\setbeamercolor{normal text}{fg=Ink}
\setbeamertemplate{navigation symbols}{}
\setbeamertemplate{headline}{}
\setbeamertemplate{footline}{}
\setbeamersize{text margin left=0pt,text margin right=0pt}
\hypersetup{colorlinks=true,linkcolor=Classical,urlcolor=Classical,pdfpagemode=UseOutlines,pdfstartview=Fit,
 pdftitle={Borealis: serie temporali ambientali},pdfauthor={Marzio De Corato},
 pdfsubject={Modulo M09: calendario, dipendenza, trasformazioni e stazionarietà, Revisione completa: diagrammi TikZ, formule e codice R commentato}}
\bookmarksetup{open,openlevel=0,depth=2}
\newsavebox{\BodyBox}
\newsavebox{\SideBox}
\newsavebox{\RecallBox}
\newsavebox{\TitleBox}
\newcommand{\MeasureBox}[2]{\typeout{REVIEWBOX:#1:\the\wd#2:\the\ht#2:\the\dp#2}}

% Revisione completa: diagrammi TikZ, formalismo e box R.
\usepackage[most]{tcolorbox}
\usepackage{listings}
\usetikzlibrary{fit,backgrounds,matrix,decorations.pathreplacing}
\definecolor{GeoA}{HTML}{1D7890}
\definecolor{GeoB}{HTML}{C57E42}
\definecolor{GeoC}{HTML}{7363AA}
\definecolor{GeoTint}{HTML}{F0F4F3}
\lstdefinestyle{GeoR}{language=R,basicstyle=\ttfamily\fontsize{9.2}{12.2}\selectfont,
 keywordstyle=\color{Quantum}\bfseries,commentstyle=\color{Classical},stringstyle=\color{GeoB},
 numbers=left,numberstyle=\tiny\color{Muted},numbersep=4pt,showstringspaces=false,
 columns=fullflexible,keepspaces=true,breaklines=true,xleftmargin=1.15em,
 frame=none,aboveskip=0pt,belowskip=0pt,tabsize=2}
\newsavebox{\PilotLeftBox}
\newsavebox{\PilotRightBox}
\newcommand{\PLabel}[1]{\par{\color{Classical}\bfseries\fontsize{9}{11}\selectfont\MakeUppercase{#1}}\par\vspace{1.8mm}}
\newcommand{\PBody}[1]{\par{\fontsize{11.3}{14.8}\selectfont #1\par}}
\newcommand{\PSmall}[1]{\par{\fontsize{9.8}{13.1}\selectfont #1\par}}
\newcommand{\PCode}[2]{%
\begin{tcolorbox}[enhanced,arc=1.8mm,boxrule=.65pt,
  colback=white,colframe=Classical!75!Navy,colbacktitle=Navy,
  coltitle=white,fonttitle=\bfseries\fontsize{9.6}{12}\selectfont,
  title={#1},left=1.7mm,right=1.3mm,top=1.5mm,bottom=1.3mm,
  before skip=1mm,after skip=1mm]
\lstinputlisting[style=GeoR]{Codice/#2}
\end{tcolorbox}}
\newcommand{\PInfo}[1]{%
\begin{tcolorbox}[enhanced,arc=1.6mm,boxrule=.55pt,colback=GeoTint,
 colframe=Rule!80!Classical,left=2.3mm,right=2.3mm,top=2mm,bottom=2mm,
 before skip=1mm,after skip=1mm]
\PSmall{#1}\end{tcolorbox}}
\tikzset{
 pFlow/.style={draw=Classical!45,fill=white,rounded corners=2mm,
  line width=.7pt,align=center,minimum height=1.1cm,font=\sffamily\fontsize{10}{12}\selectfont},
 pEdge/.style={-{Latex[length=2.3mm]},draw=Classical,line width=.95pt},
 pGroup/.style={draw=Rule,fill=white,rounded corners=2mm,line width=.65pt},
 pDotA/.style={circle,fill=GeoA,inner sep=1.25pt},
 pDotB/.style={circle,fill=GeoB,inner sep=1.25pt},
 pDotC/.style={circle,fill=GeoC,inner sep=1.25pt}
}
\newcommand{\PilotChrome}[4]{%
  \node[anchor=north west,inner sep=0,text width=17.7cm] at ([xshift=.95cm,yshift=-.55cm]current page.north west) {\color{Navy}\bfseries\fontsize{20.6}{24.5}\selectfont #1};
  \node[anchor=north east,inner sep=0] at ([xshift=19.05cm,yshift=-.14cm]current page.north west) {\color{Classical}\fontsize{8}{10}\selectfont\hyperlink{slide2}{#2}};
  \draw[Classical,line width=1.2pt] ([xshift=.95cm,yshift=-1.7cm]current page.north west)--([xshift=2.35cm,yshift=-1.7cm]current page.north west);
  \draw[Rule,line width=.65pt] ([xshift=.95cm,yshift=-10.36cm]current page.north west)--([xshift=19.05cm,yshift=-10.36cm]current page.north west);
  \node[anchor=north west,inner sep=0] at ([xshift=.96cm,yshift=-10.55cm]current page.north west) {\color{Classical}\bfseries\fontsize{9.1}{12}\selectfont DA RICORDARE};
  \node[anchor=north west,inner sep=0,text width=14.3cm] at ([xshift=4.2cm,yshift=-10.53cm]current page.north west) {\color{Ink}\fontsize{10.3}{13.8}\selectfont #4};
  \draw[Rule,line width=.35pt] ([xshift=.95cm,yshift=-11.92cm]current page.north west)--([xshift=19.05cm,yshift=-11.92cm]current page.north west);
  \node[anchor=north west,inner sep=0] at ([xshift=.95cm,yshift=-12.06cm]current page.north west) {\color{Muted}\fontsize{7.7}{9.2}\selectfont Borealis: serie temporali ambientali};
  \node[anchor=north east,inner sep=0] at ([xshift=19.05cm,yshift=-12.06cm]current page.north west) {\color{Muted}\fontsize{7.7}{9.2}\selectfont #3/106};
}


% Stili comuni per i diagrammi scientifici della revisione completa.
\tikzset{
 pAxis/.style={-{Latex[length=2mm]},draw=Muted,line width=.65pt},
 pNote/.style={font=\sffamily\fontsize{9}{11.5}\selectfont,text=Muted,align=center},
 pLabel/.style={font=\sffamily\fontsize{10}{12.5}\selectfont,text=Ink,align=center}
}
\newsavebox{\FullWideBox}
\setlength{\emergencystretch}{1em}




\newcommand{\PLab}[2]{\par{\color{Classical}\bfseries\fontsize{9}{11}\selectfont
\href{#1}{\MakeUppercase{#2}}}\par\vspace{1.8mm}}
